% !TEX root = ../main.tex

\chapter{简介}

\section{一级节标题}

\subsection{二级节标题}

\subsubsection{三级节标题}

本模板 \pkg{ustcthesis} 是中国科学技术大学本科生和研究生学位论文的 \LaTeX{}
模板， 按照《中国科学技术大学研究生学位论文撰写手册》（最近在修订中，以下简称《撰写手册》）和
《\href{https://www.teach.ustc.edu.cn/?attachment_id=13867}
{中国科学技术大学本科毕业论文（设计）格式}》的要求编写。

Lorem ipsum dolor sit amet, consectetur adipiscing elit, sed do eiusmod tempor
incididunt ut labore et dolore magna aliqua.
Ut enim ad minim veniam, quis nostrud exercitation ullamco laboris nisi ut
aliquip ex ea commodo consequat.
Duis aute irure dolor in reprehenderit in voluptate velit esse cillum dolore eu
fugiat nulla pariatur.
Excepteur sint occaecat cupidatat non proident, sunt in culpa qui officia
deserunt mollit anim id est laborum.


\section{题名页}

题名页主要包括以下内容：

A) 分类号：需同时标注中图分类号与UDC分类号两种。
选择分类号时，应以论文内容主题所涉及的科学领域为依据，
可同时参考自身专业。一般情况下选至二级或三级类目即可，
也可选至更精确的类目号。
若不确定具体分类，可咨询学校图书馆。

a) 中图分类号：采用《中国图书馆分类法》最新版本标注，
可参考：\url{https://www.clcindex.com/category}。
字符总数对应类目的级数，例如：论文中涉及火星，专业为地球物理学，
二级分类号为P1（天文学），三级分类号为P18（太阳系），
或者更详细至P185.3（火星）。

b) UDC分类号：采用《国际十进分类法》（Universal Decimal
Classification） 最新版本标注。
可参考：\url{https://udcsummary.info/php/index.php?lang=chi}。
数字总数对应类目的级数，例如：论文中涉及行星地质学，专业为地球物理学，
二级分类号为55（地球科学、地质学），三级分类号为550（地质学附属学科），
四级分类号为550.3（地球物理学）。



\section{脚注}

Lorem ipsum dolor sit amet, consectetur adipiscing elit, sed do eiusmod tempor
incididunt ut labore et dolore magna aliqua.
\footnote{Ut enim ad minim veniam, quis nostrud exercitation ullamco laboris
  nisi ut aliquip ex ea commodo consequat.
  Duis aute irure dolor in reprehenderit in voluptate velit esse cillum dolore
  eu fugiat nulla pariatur.}
